% Generated draft: requires mathematical and source review.
\documentclass[11pt]{article}
\usepackage[margin=25mm]{geometry}
\usepackage{fontspec}
\setmainfont{Latin Modern Roman}
\usepackage{amsmath,amssymb,mathtools,unicode-math}
\setmathfont{Latin Modern Math}
\usepackage{enumitem,needspace,xurl,hyperref}
\setlength{\parindent}{0pt}
\setlength{\parskip}{5pt}
\setlength{\emergencystretch}{3em}
\newcommand{\N}{\mathbb{N}}
\newcommand{\Z}{\mathbb{Z}}
\newcommand{\Q}{\mathbb{Q}}
\newcommand{\R}{\mathbb{R}}
\newcommand{\C}{\mathbb{C}}
\newcommand{\F}{\mathbb{F}}
\newcommand{\A}{\mathbb{A}}
\newcommand{\PP}{\mathbb{P}}
\newcommand{\E}{\mathbb{E}}
\newcommand{\eps}{\varepsilon}
\newcommand{\dd}{\,\mathrm{d}}
\begin{document}
\section*{UnsolvedMath ID: 800731}
\textbf{Source problem code:} catalogue entry
\subsection{Definitions and mathematical statement}
Let \(X\) be a topological space. The supplied excerpt uses \(S(X)\) without defining its elements or any topology on it. For a subgroup \(H\leq S(X)\), write \(h(x)\) for its indicated action on \(X\). The action is \emph{transitive} if
\[
\forall x,y\in X\quad \exists h\in H\quad h(x)=y.
\]
Following the supplied definition, \(H\) is \emph{distal} if
\[
\forall x_1,x_2\in X\text{ with }x_1\ne x_2,\quad
\forall x\in X,\quad
\exists U\text{ a neighborhood of }x,\quad
\forall h\in H,\quad
\{h(x_1),h(x_2)\}\not\subseteq U.
\]
Here a neighborhood of \(x\) means a set containing an open set that contains \(x\). Thus the chosen neighborhood may depend on \(x_1,x_2,x\), but must work for every \(h\in H\).

The question is: for every compact space \(X\) and every distal subgroup \(H\leq S(X)\) acting transitively on \(X\), must \(H\) be left bounded?

The terms \(S(X)\) and \emph{left bounded} must retain their intended source meanings, which are not specified in the supplied excerpt. In particular, no topology or other structure defining left boundedness is supplied, so the conclusion cannot be expanded into a precise covering condition from this evidence alone. Compactness is retained as stated; the excerpt does not explicitly specify separation axioms or exclude empty or singleton spaces.

As background, the supplied excerpt reports that \([\mathrm{Pr1},\text{ Theorem }4]\) proves that a left bounded subgroup \(H\subset S(X)\) is distal when it acts transitively on \(X\). This is a reported source claim, with the same unresolved terminology and an incomplete bibliographic reference.
\subsection{Short English statement}
For a compact space, suppose a subgroup acting transitively satisfies the stated distal condition: each distinct pair stays uniformly out of some neighborhood of every proposed common limit point under all group elements. Must that subgroup be left bounded? The supplied excerpt does not define \(S(X)\) or left boundedness, so those meanings require source clarification.
\subsection{Sources}
\begin{itemize}
\item Supplied local UnsolvedMath registry record: permanent ID \(800731\), source code \texttt{catalogue entry}, title \emph{Open Problems in Topological Algebra --- Problem 5 (Protasov)}. The registry labels the problem \emph{open}; that label is supplied metadata, not an independently verified status assessment.
\item Banakh and Choban, \emph{Some Open Problems in Topological Algebra} (2012), Problem 5. Supplied source URL: \url{https://arxiv.org/abs/1202.1619}. The record lists Banakh and Choban as proposers and gives Protasov as attribution/context.
\item Supplied provenance: accessed 2026-10-16; extraction method \texttt{arxiv\_latex}. These are supplied provenance details, not a claim of independent source consultation.
\item The excerpt cites \([\mathrm{Pr1},\text{ Theorem }4]\) for the reported implication from transitive left boundedness to distality. No full bibliographic entry for \([\mathrm{Pr1}]\) is supplied.
\end{itemize}
\end{document}
